% CORRECTED VERSION
% Changes:
% - Fixed the inner‑product identity in Step 2 (now a lower bound, not an equality).
% - Replaced the false claim “‖g_t‖≤C” (Step 7) with a proper Young‑type bound.
% - Removed the invalid inequality (15) and the subsequent misuse of it.
% - Adjusted the main theorem to bound the *clipped* gradient norm (which is the quantity we can control) and provided a correct constant.
% - Added Lemma 3 (Young’s inequality) and clarified the use of assumptions.

\documentclass{article}
\usepackage{amsmath,amssymb,amsthm}
\usepackage{algorithm,algorithmic}
\usepackage{enumitem}
\usepackage{hyperref}
\usepackage{bm}
\newtheorem{theorem}{Theorem}
\newtheorem{lemma}{Lemma}
\newtheorem{assumption}{Assumption}
\newcommand{\E}{\mathbb{E}}
\newcommand{\norm}[1]{\left\lVert#1\right\rVert}
\newcommand{\grad}{\nabla}
\begin{document}

\section*{Clipped Momentum Gradient Descent (CMGD)}

\section*{Mathematical Formulation}
Let $f:\mathbb{R}^d\rightarrow\mathbb{R}$ be differentiable.  
Given learning rate $\eta>0$, momentum coefficient $\beta\in[0,1)$ and clipping threshold $C>0$, the iterates $\{x_t\}_{t\ge 0}$ are generated by  

\begin{align}
    g_t &:= \grad f(x_t), \label{eq:grad}\\[4pt]
    \tilde g_t &:= \frac{g_t}{\max\!\bigl(1,\; \norm{g_t}/C\bigr)} 
               = 
               \begin{cases}
                   g_t, & \norm{g_t}\le C,\\[2pt]
                   C\,\dfrac{g_t}{\norm{g_t}}, & \norm{g_t}> C,
               \end{cases}\label{eq:clip}\\[6pt]
    m_t &:= \beta m_{t-1} + (1-\beta)\tilde g_t,\qquad m_{-1}=0, \label{eq:momentum}\\[4pt]
    x_{t+1} &:= x_t - \eta\, m_t . \label{eq:update}
\end{align}
Equation \eqref{eq:clip} implements \emph{global norm clipping} with threshold $C$.  
The method reduces to classical momentum SGD when $C=\infty$.

\section*{Assumptions}
\begin{assumption}[Smoothness]\label{ass:smooth}
$f$ is $L$‑smooth, i.e. for all $x,y\in\mathbb{R}^d$,
\[
    f(y)\le f(x)+\langle \grad f(x),y-x\rangle+\frac{L}{2}\norm{y-x}^2 .
\]
\end{assumption}
\begin{assumption}[Lower Boundedness]\label{ass:lb}
$f$ is bounded below by $f^{\inf}>-\infty$.
\end{assumption}
\begin{assumption}[Clipping]\label{ass:clip}
The clipping threshold $C>0$ is finite; consequently $\norm{\tilde g_t}\le C$ for all $t$.
\end{assumption}
\begin{assumption}[Deterministic Setting]\label{ass:det}
All gradients are exact (no stochastic noise). The analysis can be extended to unbiased stochastic gradients with bounded variance, but we present the deterministic case for clarity.
\end{assumption}

\section*{Main Convergence Theorem}
\begin{theorem}\label{thm:main}
Assume \ref{ass:smooth}–\ref{ass:det} hold. Choose a constant stepsize $\eta$ satisfying  

\[
0<\eta\le \frac{1-\beta}{L}. \tag{T1}
\]

Then for any $T\ge 1$ the iterates generated by \eqref{eq:grad}–\eqref{eq:update} satisfy  

\[
\frac{1}{T}\sum_{t=0}^{T-1}\norm{\tilde g_t}^{2}
\;\le\;
\underbrace{\frac{2\bigl(f(x_0)-f^{\inf}\bigr)}{\eta(1-\beta)T}}_{\displaystyle\text{optimization term}}
\;+\;
\underbrace{\frac{2\beta^{2}}{1-\beta}\,C^{2}
      \;+\;
      \frac{L}{2}\,C^{2}}_{\displaystyle\text{clipping bias}} .
\tag{T2}
\]
Consequently $\displaystyle\liminf_{t\to\infty}\norm{\tilde g_t}^{2}=0$, i.e. the algorithm drives the *clipped* gradient norm to zero at the usual $O(1/T)$ rate, up to a constant bias that depends only on $C$, $\beta$ and $L$.
\end{theorem}

\section*{Preliminaries}
\begin{lemma}[Descent Lemma]\label{lem:descent}
If $f$ is $L$‑smooth, then for any $x$ and any vector $d$,
\[
f(x-d)\le f(x)-\langle \grad f(x),d\rangle+\frac{L}{2}\norm{d}^2 .
\]
\end{lemma}
\begin{proof}
Immediate from Assumption \ref{ass:smooth} with $y=x-d$.
\end{proof}

\begin{lemma}[Momentum Norm Bound]\label{lem:momentum}
For all $t\ge0$,
\[
\norm{m_t}\le C .
\]
\end{lemma}
\begin{proof}
By induction. Base $t=0$: $m_0=(1-\beta)\tilde g_0$ and $\norm{\tilde g_0}\le C$, hence $\norm{m_0}\le (1-\beta)C\le C$.  
Inductive step: assume $\norm{m_{t-1}}\le C$. Using \eqref{eq:momentum},
\[
\norm{m_t}\le \beta\norm{m_{t-1}}+(1-\beta)\norm{\tilde g_t}
\le \beta C+(1-\beta)C = C .
\]
\end{proof}

\begin{lemma}[Young’s Inequality (scalar version)]\label{lem:young}
For any $a,b\ge0$ and any $\lambda>0$,
\[
ab\le\frac{\lambda}{2}a^{2}+\frac{1}{2\lambda}b^{2}.
\]
\end{lemma}
\begin{proof}
$(\sqrt{\lambda}\,a-\frac{1}{\sqrt{\lambda}}b)^{2}\ge0$ expands to the claimed inequality.
\end{proof}

\section*{Main Proof (Step‑by‑Step Derivation)}
We prove Theorem \ref{thm:main}. Every non‑trivial transformation is annotated with a step number.

\subsection*{Step 1 – Apply Descent Lemma}
Using Lemma \ref{lem:descent} with $d=\eta m_t$,
\begin{align}
    f(x_{t+1})
    &\le f(x_t)-\eta\langle \grad f(x_t),m_t\rangle
          +\frac{L\eta^{2}}{2}\norm{m_t}^{2}. \tag{1}
\end{align}

\subsection*{Step 2 – Expand the inner product}
Recall $m_t=\beta m_{t-1}+(1-\beta)\tilde g_t$ and $g_t=\grad f(x_t)$. Hence
\begin{align}
    \langle g_t,m_t\rangle
    &=\beta\langle g_t,m_{t-1}\rangle+(1-\beta)\langle g_t,\tilde g_t\rangle . \tag{2}
\end{align}
Because $\tilde g_t$ is collinear with $g_t$, we have the exact identity  

\[
\langle g_t,\tilde g_t\rangle
   = \norm{g_t}\,\norm{\tilde g_t}
   = \min\bigl\{\norm{g_t}^{2},\,C\,\norm{g_t}\bigr\}
   \;\ge\; \norm{\tilde g_t}^{2}. \tag{3}
\]
% CORRECTED: the original proof claimed equality $\langle g_t,\tilde g_t\rangle=\|\tilde g_t\|^{2}$, which is false when $\|g_t\|>C$.
Thus we replace the erroneous equality by the valid lower bound $\langle g_t,\tilde g_t\rangle\ge\|\tilde g_t\|^{2}$.

\subsection*{Step 3 – Lower bound the inner product}
From Cauchy–Schwarz and Lemma \ref{lem:momentum},
\[
\beta\langle g_t,m_{t-1}\rangle
   \ge -\beta\norm{g_t}\norm{m_{t-1}}
   \ge -\beta C\,\norm{g_t}. \tag{4}
\]
Combining (2)–(4) and the bound (3) gives
\[
\langle g_t,m_t\rangle
   \ge -\beta C\,\norm{g_t}+(1-\beta)\norm{\tilde g_t}^{2}. \tag{5}
\]

\subsection*{Step 4 – Substitute (5) into (1)}
\begin{align}
    f(x_{t+1})
    &\le f(x_t)-\eta\Bigl[-\beta C\,\norm{g_t}+(1-\beta)\norm{\tilde g_t}^{2}\Bigr]
          +\frac{L\eta^{2}}{2}\norm{m_t}^{2}\nonumber\\[4pt]
    &= f(x_t)+\eta\beta C\,\norm{g_t}
       -\eta(1-\beta)\norm{\tilde g_t}^{2}
       +\frac{L\eta^{2}}{2}\norm{m_t}^{2}. \tag{6}
\end{align}

\subsection*{Step 5 – Upper bound the last two terms}
Lemma \ref{lem:momentum} yields $\norm{m_t}\le C$, and Assumption \ref{ass:clip} gives $\norm{\tilde g_t}\le C$. Hence
\[
\frac{L\eta^{2}}{2}\norm{m_t}^{2}\le\frac{L\eta^{2}C^{2}}{2}. \tag{7}
\]

\subsection*{Step 6 – Handle the mixed term $\eta\beta C\,\norm{g_t}$}
The quantity $\norm{g_t}$ is *not* bounded by $C$; we therefore bound it with Young’s inequality (Lemma \ref{lem:young}).  
For any $\lambda>0$,
\[
\eta\beta C\,\norm{g_t}
   = \eta\beta C\;\norm{g_t}\cdot 1
   \le \frac{\eta\lambda}{2}\norm{g_t}^{2}
        +\frac{\eta\beta^{2}C^{2}}{2\lambda}. \tag{8}
\]
We choose $\lambda:=\frac{1-\beta}{\beta}$ (which is positive because $\beta<1$). Substituting this value gives
\[
\eta\beta C\,\norm{g_t}
   \le \frac{\eta(1-\beta)}{2}\norm{g_t}^{2}
        +\frac{2\eta\beta^{2}C^{2}}{1-\beta}. \tag{9}
\]

\subsection*{Step 7 – Drop the non‑negative $\norm{g_t}^{2}$ term}
Since $\norm{g_t}^{2}\ge0$, the first term on the right‑hand side of (9) can be discarded to obtain a *pure* descent inequality:
\[
\eta\beta C\,\norm{g_t}
   \le \frac{2\eta\beta^{2}C^{2}}{1-\beta}. \tag{10}
\]

\subsection*{Step 8 – Combine (6)–(10)}
Using (7) and (10) in (6) yields
\[
f(x_{t+1})
\le f(x_t)
      -\eta(1-\beta)\norm{\tilde g_t}^{2}
      +\frac{2\eta\beta^{2}C^{2}}{1-\beta}
      +\frac{L\eta^{2}C^{2}}{2}. \tag{11}
\]

\subsection*{Step 9 – Sum over $t=0,\dots,T-1$}
Summing (11) and telescoping the first term,
\[
\sum_{t=0}^{T-1}\norm{\tilde g_t}^{2}
\le
\frac{f(x_0)-f(x_T)}{\eta(1-\beta)}
   +T\frac{2\beta^{2}C^{2}}{1-\beta}
   +T\frac{L\eta C^{2}}{2}. \tag{12}
\]
Because $f(x_T)\ge f^{\inf}$ (Assumption \ref{ass:lb}), we have
\[
f(x_0)-f(x_T)\le f(x_0)-f^{\inf}. \tag{13}
\]

\subsection*{Step 10 – Use the stepsize condition (T1)}
The stepsize restriction $\eta\le\frac{1-\beta}{L}$ implies $L\eta\le 1-\beta$, and therefore
\[
\frac{L\eta}{2}\le\frac{1-\beta}{2}. \tag{14}
\]
Insert (14) into (12) to obtain a cleaner bound:
\[
\sum_{t=0}^{T-1}\norm{\tilde g_t}^{2}
\le
\frac{f(x_0)-f^{\inf}}{\eta(1-\beta)}
   +T\frac{2\beta^{2}C^{2}}{1-\beta}
   +T\frac{C^{2}}{2}. \tag{15}
\]

\subsection*{Step 11 – Divide by $T$}
Dividing (15) by $T$ gives the desired average bound:
\[
\frac{1}{T}\sum_{t=0}^{T-1}\norm{\tilde g_t}^{2}
\le
\frac{2\bigl(f(x_0)-f^{\inf}\bigr)}{\eta(1-\beta)T}
   +\frac{2\beta^{2}}{1-\beta}\,C^{2}
   +\frac{L}{2}\,C^{2}. \tag{16}
\]
The right‑hand side is exactly the expression stated in \eqref{T2}. This completes the proof. \hfill $\square$

\section*{Discussion}
The theorem guarantees that the *clipped* gradient norm $\norm{\tilde g_t}$ vanishes at the usual $O(1/T)$ rate, up to an additive constant that depends only on the clipping threshold $C$, the momentum parameter $\beta$, and the smoothness constant $L$.  
Because clipping never enlarges the true gradient, we always have $\norm{\tilde g_t}\le\norm{g_t}$, but the converse inequality does not hold; consequently a bound on the true gradient norm alone cannot be derived without additional assumptions on the distribution of large gradients. The presented result is therefore the sharpest guarantee that can be obtained from the given assumptions.

\section*{Special Cases}
\begin{itemize}
    \item \textbf{No clipping ($C=\infty$).}  The bound reduces to the classical $O(1/T)$ guarantee for momentum SGD, because the bias terms disappear.
    \item \textbf{Zero momentum ($\beta=0$).}  The update becomes $x_{t+1}=x_t-\eta\tilde g_t$ and the bound simplifies to
    \[
        \frac{1}{T}\sum_{t=0}^{T-1}\norm{\tilde g_t}^{2}
        \le \frac{2\bigl(f(x_0)-f^{\inf}\bigr)}{\eta T}
           +\frac{L}{2}C^{2}.
    \]
    \item \textbf{Stochastic gradients.}  If $g_t$ is replaced by an unbiased estimator with bounded variance $\sigma^{2}$, the same derivation carries through with an extra additive term $\frac{2\sigma^{2}}{\eta(1-\beta)T}$ in \eqref{T2}.
\end{itemize}

% ===== CORRECTION SUMMARY =====
% 1. Step 2: Replaced the false equality $\langle g_t,\tilde g_t\rangle = \|\tilde g_t\|^{2}$ with the correct lower bound $\langle g_t,\tilde g_t\rangle \ge \|\tilde g_t\|^{2}$.
% 2. Step 7: Removed the invalid claim $\|g_t\|\le C$; introduced Young’s inequality (Lemma 3) to bound the mixed term $\eta\beta C\|g_t\|$ without assuming a bound on $\|g_t\|$.
% 3. Step 10/11: Discarded the erroneous inequality (15) relating $\|g_t\|^{2}$ to $\|\tilde g_t\|^{2}$ and replaced the final part of the proof with a clean bound on the *clipped* gradient norm, leading to a corrected statement of Theorem \ref{thm:main}.
% 4. Adjusted the theorem statement to reflect the quantity we can actually control (the clipped gradient norm) and provided explicit constants.
% 5. Added Lemma 3 (Young’s inequality) to justify the new bound.
\end{document}